\subsection{Unitary group and orthogonal group.}

Let \(\Phi\) be a Hermitian form on E; the automorphisms of the A-module E that leave \(\Phi\) invariant are called unitary automorphisms (or unitary transformations) relative to \(\Phi\) , and their group is called the unitary group associated with \(\Phi\) ; it is denoted \(\mathbf{U}(\Phi)\) . Given a quadratic form \(Q \neq 0\) on E, the automorphisms of the A-module E which leave Q invariant are called orthogonal automorphisms (or orthogonal transformations) relative to Q; their group is called the orthogonal group associated with Q; it is denoted by \(\mathbf{O}(Q)\) .

Every orthogonal transformation for a quadratic form Q is unitary for the bilinear form associated with Q. The converse is true when the scalar 2 is not equal to 0 or a zero divisor in A (§ 3, no. 4, (13)), for example if A is a field of characteristic ≠ 2.

Consider, in particular, on the module \(E = A^{n}\) the Hermitian form \(\Phi_{0}\) whose matrix with respect to the canonical basis \((e_{i})\) of E is the identity matrix \(I_{n}\) . The unitary automorphisms associated with \(\Phi_{0}\) are simply called unitary automorphisms (or transformations) in n variables; their group is called the unitary group in n variables and is sometimes denoted \(\mathbf{U}(n, \mathbf{A})\) or \(\mathbf{U}_{n}(\mathbf{A})\) . The matrix U of a unitary automorphism with respect to \((e_{i})\) is called a unitary matrix. Such a matrix is invertible, and satisfies, according to formula (48) of § 1, no. 10, the relation

\[
{ } ^ { t } U . \overline { { U } } = I _ { n } ;\tag{4}
\]

Conversely, if A is a commutative ring or a field, a matrix U satisfying (4) is invertible, and is then unitary.

When J is the identity and 2 is neither equal to 0 nor a divisor of 0 in A, one uses the terms orthogonal group in n variables, orthogonal automorphism (or orthogonal transformation) in n variables, and orthogonal matrix instead of the preceding terms, and one writes \(\mathbf{O}(n,\mathrm{A})\) (resp. \(\mathbf{O}_n(\mathrm{A})\) ) instead of \(\mathbf{U}(n,\mathrm{A})\) (resp. \(\mathbf{U}_n(\mathrm{A})\) ). Relation (4) is then written

\[
{ } ^ { t } U . U = I _ { n }\tag{5}
\]

and, since A is commutative, it is a necessary and sufficient condition for U to be an orthogonal matrix.

PROPOSITION 3. --- Suppose that A is a commutative field and that E is of finite dimension \textgreater{} 0. Let Φ be a nondegenerate Hermitian form on E. The map u → det u is a homomorphism from the unitary group U(Φ) associated with Φ onto the multiplicative subgroup H of A consisting of the elements ρ such that ρρ̄ = 1 (a subgroup reduced to \{1, -1\} when J is the identity).

Indeed, let \(u\) be an element of \(\mathbf{U}(\Phi)\) , \(U\) its matrix with respect to a basis of E, and \(R\) the matrix of \(\Phi\) with respect to this basis. The relation \(R = {}^t U.R.\overline{U}\) ( \(\S 1\) , no. 10, formula (48)) shows that one has \((\det U)(\det \overline{U}) = 1\) since \(R\) is invertible; whence \((\det u)(\overline{\det u}) = 1\) . The homomorphism \(u \to \det u\) maps \(\mathbf{U}(\Phi)\) onto H. Indeed, when A is of characteristic 2 and J is the identity, H is reduced to the element 1. Otherwise there exists an orthogonal basis \((e_i)(i = 1,\ldots,n)\) of E (th. 1); for every \(\rho \in A\) such that \(\rho \bar{\rho} = 1\) , let \(u\) be the automorphism of E defined by \(u(e_1) = \rho e_1\) and \(u(e_i) = e_i\) for \(i = 2,\ldots,n\) ; then \(u\) is unitary and det \(u = \rho\) , whence the proposition.

Under the conditions of prop. 3, the kernel of the homomorphism \(u \to \det u\) is a normal subgroup of \(\mathbf{U}(\Phi)\) , which is called the special unitary group associated with \(\Phi\) ; it is sometimes denoted \(\mathbf{SU}(\Phi)\) .

When J is the identity and A is not of characteristic 2, this group is also called the special orthogonal group associated with \(\Phi\) (or with the quadratic form \(Q(x) = \Phi(x, x)\) ) and is sometimes denoted SO(Q).

If E = A \(^{n}\) and if Φ is the form whose matrix with respect to the canonical basis of E is the identity matrix, one uses the notations SU(n, A) or SU \(_{n}\) (A) and SO(n, A) or SO \(_{n}\) (A).

